% plik wygenerowany przez stats/tables.py; zmiany wprowadzać w skrypcie, nie tutaj
% F1 każdej klasy dla RF (macierz 6 x 6 zsumowana po M1-M6, rf_class_f1.csv);
% klasy c2_beaconing nie ma w treningu, więc jej F1 = 0 z definicji
\begin{tabular}{lrcccc}
\toprule
Klasa & Przepływy & \begin{tabular}[b]{@{}c@{}}CIC \\ {}S1\end{tabular} & \begin{tabular}[b]{@{}c@{}}CIC \\ {}S2\end{tabular} & \begin{tabular}[b]{@{}c@{}}CIC + GeNIS \\ {}S1\end{tabular} & \begin{tabular}[b]{@{}c@{}}CIC + GeNIS \\ {}S2\end{tabular} \\
\midrule
ruch normalny & 25\,602 & 0,684 & 0,985 & 0,706 & 0,992 \\
A1 skan portów & 122\,947 & 1,000 & 1,000 & 1,000 & 1,000 \\
A2 brute-force SSH & 414 & 0,210 & 0,886 & 0,165 & 0,952 \\
A3 wolny DoS & 5097 & 0,149 & 0,986 & 0,188 & 0,994 \\
A4 HTTP flood & 49\,156 & 0,000 & 0,994 & 0,245 & 0,996 \\
Z1 beaconing C2 & 168 & 0,000 & 0,000 & 0,000 & 0,000 \\
\midrule
macro-F1 &  & 0,341 & 0,808 & 0,384 & 0,822 \\
macro-F1 bez Z1 &  & 0,409 & 0,970 & 0,461 & 0,987 \\
\bottomrule
\end{tabular}
